\documentclass[../../main.tex]{subfiles}
\begin{document}

\begin{center}
  {\fontsize{48pt}{58pt}\selectfont \textbf{2009}\normalsize\textbf{年度}}
\end{center}

\begin{table}[htb]
  \centering
  \caption{出題テーマと難易度}
  \begin{tabular}{|c|p{9cm}|c|}
    \hline
    \textbf{問題番号} & \textbf{テーマ} & \textbf{難易度} \\
    \hline
    第1問 & 放物線の直交する2接線と面積の最小値 & \\
    \hline
    第2問 & 1次変換で不変な直線の条件 & \\
    \hline
    第3問 & 整数問題（格子点の個数） & \\
    \hline
    第4問 & 空間図形の回転体の体積 & \\
    \hline
  \end{tabular}
\end{table}

\begin{table}[htb]
  \centering
  \caption{解答}
  \renewcommand{\arraystretch}{1.4}
  \begin{tabular}{|c|c|p{9cm}|}
    \hline
    \textbf{問題番号} & \textbf{小問} & \textbf{答え} \\
    \hline
    第1問 & - & $S$の最小値 $\dfrac{1}{3}$ \\
    \hline
    第2問 & - & $a=2,\ \dfrac{3}{2}$ \\
    \hline
    第3問 & - & $N(2N-1)$ \\
    \hline
    第4問 & (1) & $x+y=t$，$z=0$ \\
    \cline{2-3}
    & (2) & $V=\dfrac{2\sqrt{3}}{45}\pi$ \\
    \hline
  \end{tabular}
\end{table}

\paragraph{試験形態}
4問．試験時間は要確認．

\end{document}
